% Figure: phase and group velocity above cutoff in a rectangular waveguide,
% normalised to c, versus normalised frequency f/f_c. Phase velocity exceeds c,
% group velocity falls below c, both approach c far above cutoff.
\begin{figure}[htbp]
\centering
\begin{tikzpicture}
\begin{axis}[
  width=11cm, height=6cm,
  xlabel={frequency $f/f_c$},
  ylabel={velocity $/\,c$},
  xmin=1, xmax=4, ymin=0, ymax=3,
  domain=1.02:4, samples=200,
  axis lines=left,
  legend pos=north east,
  legend cell align=left,
  enlargelimits=false,
]
  % phase velocity v_p/c = 1/sqrt(1 - (1/x)^2)
  \addplot[engblue, very thick] { 1/sqrt(1 - (1/x)^2) };
  \addlegendentry{$v_p/c$}
  % group velocity v_g/c = sqrt(1 - (1/x)^2)
  \addplot[engred, very thick] { sqrt(1 - (1/x)^2) };
  \addlegendentry{$v_g/c$}
  % c reference
  \addplot[enggray, dashed, thick, domain=1:4] {1};
  \addlegendentry{$c$}
\end{axis}
\end{tikzpicture}
\caption[Phase and group velocity in a waveguide above cutoff]{Phase
velocity $v_p = c/\sqrt{1-(f_c/f)^2}$ and group velocity $v_g =
c\sqrt{1-(f_c/f)^2}$ for the $\text{TE}_{10}$ mode of a rectangular
waveguide, normalised to $c$ and plotted against $f/f_c$. Near cutoff
the phase velocity diverges and the group velocity falls to zero;
far above cutoff both approach $c$. Energy and information travel at
$v_g$; the superluminal phase velocity carries neither and violates no
relativity.}
\label{fig:vol04ch10:waveguide-dispersion}
\end{figure}
